\documentclass[11pt,a4paper]{article}
\usepackage[T1]{fontenc}
\usepackage[utf8]{inputenc}
\usepackage{lmodern,amsmath,amssymb}
\usepackage[margin=25mm]{geometry}
\usepackage{longtable,booktabs,array,calc}
\usepackage[hidelinks]{hyperref}
\hypersetup{pdftitle={What survives refinement: the pion, the Yang--Mills gap and Newton's action},pdfauthor={navstokgap research working notes}}
\newcommand{\tightlist}{\setlength{\itemsep}{0pt}\setlength{\parskip}{0pt}}
\setlength{\emergencystretch}{3em}
\setcounter{secnumdepth}{0}
\begin{document}
\hypertarget{what-survives-refinement-the-pion-the-yangmills-gap-and-newtons-action}{%
\section{What survives refinement: the pion, the Yang--Mills gap and
Newton's
action}\label{what-survives-refinement-the-pion-the-yangmills-gap-and-newtons-action}}

\textbf{Working conclusion, 2026-09-26.} The common problem is to
construct a limit while preserving a specified physical structure. In
pure \(SU(3)\) Yang--Mills it is a positive energy threshold. In chiral
QCD it is the broken global symmetry and observable spectral weight that
require a zero threshold. In Newton's comparison it is a positive action
cost of physical records while the time partition becomes arbitrarily
fine. All three require control independent of the regulator. The
quantities controlled, and the reasons for their survival, differ.

User direction of this date makes this comparison the intended frame of
the eventual paper. \textbf{The two main proof goals remain the pure
\(SU(3)\) mass gap and the necessity of a positive action scale in
Newton's comparison.} QCD with massless and with nonzero quark masses is
an orientation and a source of mechanisms, rather than a third full
construction programme. This note supplies the definitions, two
elementary limit criteria, and the plan. It proves neither main goal.

The \href{planck-gap-paper.md}{Planck paper} remains the developed
Newton component; the \href{mass-gap-position.md}{Yang--Mills position}
and \href{mass-gap-conditional-theorem.md}{conditional theorem} remain
the technical map for pure gauge theory. Their results retain their
stated hypotheses. The present note changes the destination of the
synthesis.

\hypertarget{compare-physical-quantities-with-the-regulators-separate}{%
\subsection{1. Compare physical quantities, with the regulators
separate}\label{compare-physical-quantities-with-the-regulators-separate}}

Use \(a\) for the spacetime lattice spacing, \(L\) for spatial box size,
\(m_q\) for a renormalized quark mass, and \(\varepsilon=|\pi|\) for the
largest step of a Newtonian time partition. Euclidean time extent is
taken to infinity to select zero temperature. Write \(E_*\) for an
energy threshold, \(M_*=E_*/c^2\) for a mass, and
\(\mu_*=E_*/(\hbar c)\) for an inverse correlation length. This avoids
the convention in older notes where \(m\) sometimes denotes an energy.

For an isotropic reflection-positive lattice theory, on a specified
physical channel, normalize the transfer matrix by its vacuum eigenvalue
and write

\[T_a=e^{-aH_a/(\hbar c)},\qquad
\delta_a=\frac{aE_*(a)}{\hbar c}=a\mu_*(a).\]

A continuum theory with \(0<E_*<\infty\) therefore has \(\delta_a\to0\).
A massless channel also has \(\delta_a\to0\). Their distinction is the
limit of \(\delta_a/a\), together with convergence of the states and
observables that detect it. A strong-coupling lattice gap at fixed bare
coupling does not specify that limit.

\begin{longtable}[]{@{}
  >{\raggedright\arraybackslash}p{(\columnwidth - 4\tabcolsep) * \real{0.3333}}
  >{\raggedright\arraybackslash}p{(\columnwidth - 4\tabcolsep) * \real{0.3333}}
  >{\raggedright\arraybackslash}p{(\columnwidth - 4\tabcolsep) * \real{0.3333}}@{}}
\toprule\noalign{}
\begin{minipage}[b]{\linewidth}\raggedright
Comparison
\end{minipage} & \begin{minipage}[b]{\linewidth}\raggedright
Quantity to retain in physical units
\end{minipage} & \begin{minipage}[b]{\linewidth}\raggedright
Required limiting structure
\end{minipage} \\
\midrule\noalign{}
\endhead
\bottomrule\noalign{}
\endlastfoot
Pure \(SU(3)\) & \(0<E_{\rm YM}<\infty\) & A nontrivial continuum
theory; a gap above its vacuum \\
Two-flavour QCD, \(m_q=0\) & \(E_\pi=0\) with nonzero pion spectral
weight & Conserved non-singlet axial current and a broken-symmetry
vacuum \\
Two-flavour QCD, small \(m_q>0\) & \(E_\pi(m_q)>0\), tending to zero
with \(m_q\) & Controlled explicit symmetry breaking and the same
physical scale \\
Newton comparison & \(h_*>0\) in a specified class of physical records &
A record algebra and cost surviving \(\varepsilon\to0\) \\
\end{longtable}

The first three rows use a quantum theory with \(\hbar>0\) already in
its definition. The last main goal asks how that positive action
structure is justified. The pion supplies a useful check: positive
\(\hbar\) and a dynamically generated strong-interaction scale can
coexist with zero energy gap.

\hypertarget{pure-yangmills-retain-coercivity-through-the-limit}{%
\subsection{2. Pure Yang--Mills: retain coercivity through the
limit}\label{pure-yangmills-retain-coercivity-through-the-limit}}

The target is a nontrivial four-dimensional \(SU(3)\) theory with the
Wightman or Osterwalder--Schrader properties, and physical Hamiltonian
\(H\) satisfying

\[H\Omega=0,\qquad
E_{\rm YM}:=\inf\bigl(\operatorname{spec}H\setminus\{0\}\bigr)
\in(0,\infty).\]

One chooses a vacuum sector and proves the corresponding vacuum and
clustering properties. This target requires a gap above the vacuum; it
does not prescribe every point of the spectrum above the threshold. The
official statement explicitly couples existence, axioms and the gap
(\href{https://www.claymath.org/wp-content/uploads/2022/06/yangmills.pdf}{Jaffe--Witten,
§§4--6}, passage).

For the lattice route, fix a renormalized physical scale and tune
\(g_0(a)\) toward zero. A sufficient spectral estimate has the form

\[H_{a,L}-E_{0,a,L}\ \ge\ E_-(1-P_{0,a,L}),\qquad E_->0,\]

uniformly along the trajectory in sufficiently small \(a\) and
arbitrarily large \(L\), on the gauge-invariant Hilbert space. To pass
this estimate to a continuum theory one still needs convergence of
renormalized local observables, positivity, covariance, regularity and
nontriviality. An upper estimate or a surviving finite-energy excitation
prevents all non-vacuum states from escaping to infinite energy.

A corresponding correlation route is a bound
\(C_O(t)\le K_Oe^{-E_-t/\hbar}\) for positive Euclidean time
autocorrelations of a dense family of vacuum-subtracted local vectors.
The constants must survive the cutoff limit for their chosen
renormalizations. One rapidly decaying glueball correlator establishes
only a channel statement. Reflection positivity and reconstruction are
part of the route from correlations to the physical Hamiltonian; a
Langevin relaxation rate by itself is a different object.

The existing H1/H2 route aims to deliver
\(E_-=(\hbar c/a_*)\min(\gamma_*,\gamma')\). It remains conditional on
the blocking and mixing estimates. Construction and nontriviality,
called T4 in the \href{mass-gap-obligations-lattice.md}{obligations
map}, must remain explicit as well. We choose \(L\to\infty\) at fixed
\(a\) followed by \(a\to0\) along the trajectory as the working order;
interchanging these operations requires its own uniform estimates.

\hypertarget{the-pion-preserve-the-symmetry-that-requires-a-soft-channel}{%
\subsection{3. The pion: preserve the symmetry that requires a soft
channel}\label{the-pion-preserve-the-symmetry-that-requires-a-soft-channel}}

Take \(SU(3)\) colour with two degenerate quarks, zero temperature,
vacuum angle zero, and no electromagnetism. This fixes what is meant by
the comparison with and without a quark mass. Additional flavours or
electromagnetism would change the statement.

\textbf{Massless case.} Conditional on a local relativistic continuum
theory and spontaneous breaking
\(SU(2)_L\mathbin{\times}SU(2)_R\to SU(2)_V\), the Goldstone theorem
requires three massless modes with nonzero axial-current matrix
elements. These are the pion modes. The theorem's symmetry-breaking
premise is a dynamical obligation in QCD; the theorem does not establish
it merely from the presence of fermions
(\href{https://doi.org/10.1103/PhysRev.127.965}{Goldstone--Salam--Weinberg
1962}, abstract; \href{https://arxiv.org/pdf/hep-ph/9311274}{Leutwyler
1993, §§2--3}, passage).

A possible order of limits for the per-flavour order parameter is

\[\Sigma_R=-\lim_{m_q\downarrow0}\lim_{L\to\infty}\lim_{a\to0}
\langle\bar u u\rangle_{R,a,L,m_q}>0,\]

with a positive mass selecting the vacuum and the physical scale held
fixed. The inner continuum limit is taken at fixed physical \(L,m_q\).
An exactly symmetric finite-volume state has zero symmetry-breaking
expectation at zero source; it cannot select the required vacuum by
itself. Finite-volume zero modes and the infinite-volume order parameter
must be distinguished
(\href{https://arxiv.org/pdf/0812.2797}{Damgaard--Fukaya 2009,
introduction}, passage). The displayed order is a proposed construction
order, with existence at each stage an obligation.

\textbf{Nonzero case.} At fixed small positive \(m_q\), removing \(a\)
retains explicit chiral breaking. In the conventional broken-phase
low-energy description, put \(e_q=m_qc^2\) and let \(B_E>0\) have energy
units. Then

\[E_\pi^2=2B_Ee_q[1+r(e_q)],\qquad r(e_q)\longrightarrow0
\quad(e_q\downarrow0).\]

This is the leading chiral mass relation, conditional on the broken
phase and its controlled low-energy expansion; higher orders include
chiral logarithms. In particular, \(a\to0\) at fixed \(e_q>0\) is
different from \(e_q\to0\)
(\href{https://arxiv.org/pdf/hep-ph/9609467}{Leutwyler, \emph{Light
Quark Masses}, eq. (1)}, passage). We use this as a benchmark, without
claiming a constructive proof of massive QCD or a lower bound on every
channel from this formula.

There is a sharper identity underneath the expansion. Define
\(\mu_q=m_qc/\hbar\) and \(\mu_\pi=E_\pi/(\hbar c)\), both inverse
lengths. Assume the renormalized non-singlet Ward identity in compatible
Euclidean conventions, away from contact insertions:

\[\partial_\alpha A^b_\alpha=2\mu_q P^b.\]

With pole amplitudes defined by
\(\langle0|A^b_\alpha|\pi^d(k)\rangle=i\delta^{bd}{\cal F}_\pi k_\alpha\)
and \(\langle0|P^b|\pi^d\rangle=\delta^{bd}{\cal G}_\pi\), its one-pion
matrix element reads

\[\mu_\pi^2{\cal F}_\pi=2\mu_q{\cal G}_\pi.\]

The conventional continuation fixes the displayed phases. The equation
is exact once the current identity and the pion pole exist; deriving the
leading mass formula also uses the small-mass behaviour of the residues.
A pole with nonzero residue must survive the limit. An unrenormalized
divergent pseudoscalar susceptibility is insufficient: contact terms and
ultraviolet divergences need separate control.

At finite \(a\) the regulator matters. A Ginsparg--Wilson Dirac operator
has an exact lattice non-singlet chiral symmetry
(\href{https://arxiv.org/abs/hep-lat/9802011}{Lüscher 1998}, abstract).
Wilson fermions instead require critical-mass tuning and a restored,
renormalized Ward identity
(\href{https://doi.org/10.1140/epjc/s10052-019-7287-1}{Bali et al.~2019,
§2}, passage). Exact lattice chirality alone supplies neither
spontaneous breaking nor a continuum construction.

For the final paper this comparison earns its place by testing the
mechanism. A proposed pure-gauge estimate cannot be carried unchanged to
chirally broken QCD if it gaps the pion channel. Colour gauge symmetry
and global axial flavour symmetry have different roles. Supersymmetric
cancellation of transverse zero-point energies, recorded in
\href{low-dimensional-mass-gap.md}{G07}, is a separate mechanism from
QCD's Goldstone pions. Nor does a massless pion imply a massless
coloured gluon.

\hypertarget{newton-retain-a-record-cost-while-the-polygon-error-vanishes}{%
\subsection{4. Newton: retain a record cost while the polygon error
vanishes}\label{newton-retain-a-record-cost-while-the-polygon-error-vanishes}}

Use \(M\) for the body's mass, constant force \(F\), horizontal speed
\(v\), and comparison duration \(T\). The specified Galileo comparison
is

\[q_I(t)=(vt,0),\qquad q_F(t)=(vt,Ft^2/2M),\]
\[A(T)=\frac{vFT^3}{6M},\qquad
{\cal A}(T):=\frac{3F}{v}A(T)=\frac{F^2T^3}{2M}=T\Delta E.\]

For a partition \(\pi\) of \([0,T]\), the action difference associated
with its inscribed chord segments is

\[K_\pi=\frac{F^2}{24M}\sum_j(\Delta t_j)^3,
\qquad 0\le K_\pi\le\frac{F^2T}{24M}\varepsilon^2\to0.\]

This follows by \((\Delta t_j)^3\le\varepsilon^2\Delta t_j\) and
\(\sum_j\Delta t_j=T\). In the supplied quantum lift the phase
difference is \(K_\pi/\hbar\), which also tends to zero at fixed
\(\hbar>0\). The \href{polygon-lift-phase.md}{polygon note} separately
gives an ordering phase \(\Phi_\pi\) with

\[\hbar\Phi_\pi+4K_\pi=\frac{{\cal A}(T)}3.\]

Thus one geometric action defect vanishes while the ordering comparison
retains a finite action at fixed \(T\). Reading that phase requires a
relative-phase experiment. Its derivation already assumes the Weyl
relations and positive \(\hbar\).

The proposed \(h\) gap concerns what a physical record must cost. For a
fixed class \(\mathfrak M_\varepsilon\) of calibrated marks, define

\[\kappa_\varepsilon
=\inf_{M_j\in\mathfrak M_\varepsilon}\delta_j\Delta_j,\]

where \(\delta_j\) is reading error and \(\Delta_j\) impulse
uncertainty. The class, available preparations, apparatus resources and
meaning of the uncertainty must be specified, with feasible informative
records of finite cost retained in the limit. A restricted theorem might
establish
\(\inf_{0<\varepsilon<\varepsilon_0}\kappa_\varepsilon\ge\kappa_*>0\).
The general-instrument formulation instead bounds the disturbance
functional defined in the \href{record-costs-disturbance.md}{disturbance
note}. These are distinct formulations of record cost; a single-mark
product does not by itself establish a universal protocol bound.

The current paper proves such conditional bounds using quantum
kinematics or an assumed positive mark trade-off. The main necessity
goal is to supply independently justified physical premises that force a
positive scale and its universality, without placing that scale in the
premise. Identification with \(h=2\pi\hbar\) is a further calibration
obligation. Newton's optical fits provide a candidate product
\(\Lambda p\); the \href{planck-gap-paper.md}{paper's §8} keeps its
extra indeterminacy and cross-colour universality premises visible.

Two refinement operations must stay separate: making more marks within
fixed \(T\), and asking each shrinking cell to decide the force on its
own. The Gaussian result bounds the latter decision window; it permits
marks much denser than that window. A positive action cost is compatible
with continuous time and with convergence of the polygon. It is also
compatible with zero energy gap, as the quantum free particle
illustrates.

\hypertarget{two-elementary-criteria-for-carrying-spectral-information-to-a-limit}{%
\subsection{5. Two elementary criteria for carrying spectral information
to a
limit}\label{two-elementary-criteria-for-carrying-spectral-information-to-a-limit}}

These statements are standard consequences of the spectral theorem and
weak convergence of positive measures, included to state exactly what
the comparison needs. They are written derivations, with no numerical
test and no claim of novelty.

\textbf{Proposition 1 (a surviving lower bound).} Suppose renormalized,
vacuum-subtracted local vectors \(O_n\Omega_n\) have finite positive
energy spectral measures \(\nu_n\) converging weakly to \(\nu\) in a
reconstructed limiting theory. If every \(\nu_n\) is supported in
\([E_-,\infty)\) for one \(E_->0\), then so is \(\nu\).

\emph{Proof.} Every nonnegative continuous test function compactly
supported in \([0,E_-)\) integrates to zero against every \(\nu_n\),
hence against \(\nu\). Such test functions exhaust the interval. If the
statement holds on a dense family of local vectors orthogonal to the
vacuum, the spectral projection of \(H\) onto \((0,E_-)\) is zero. A
nonzero limiting measure supplies surviving finite-energy spectral
content. \(\square\)

For positive autocorrelations \(C_n(t)=\int e^{-Et/\hbar}\,d\nu_n(E)\),
a uniform physical decay rate is another way to obtain the support
condition. The normalization of the observables and their nonzero limit
are part of the hypotheses.

\textbf{Proposition 2 (a surviving zero threshold).} Suppose instead
that \(\nu_n\Rightarrow\nu\), that \(\nu(\{0\})=0\), and that for every
\(\eta>0\) there is \(w_\eta>0\) such that, for all sufficiently large
\(n\),

\[\nu_n([0,\eta/2])\ge w_\eta.\]

Then \(\nu((0,\eta))>0\) for every \(\eta>0\), and the channel's excited
energy threshold is zero.

\emph{Proof.} The closed-set inequality for weak convergence gives

\[\nu([0,\eta/2])\ge\limsup_n\nu_n([0,\eta/2])\ge w_\eta.\]

Removing the zero atom leaves positive weight in \((0,\eta)\).
\(\square\)

The no-atom hypothesis distinguishes arbitrarily soft excitations from
an extra vacuum contribution. It belongs to a selected pure vacuum and a
connected observable. The chiral comparison suggests obtaining the
required surviving weight from an axial Ward identity and a nonzero
order parameter, with the Goldstone theorem supplying the continuum
implication.

\textbf{Why convergence of eigenvalues alone is too weak.} For \(E_0>0\)
let \(\mathrm d_x\) denote unit point mass at energy \(x\), and set

\[\nu_n=(1-e^{-n})\mathrm d_{E_0}+e^{-n}\mathrm d_{E_0/n}.\]

Every \(\nu_n\) has lowest energy \(E_0/n\to0\), yet
\(\nu_n\Rightarrow\mathrm d_{E_0}\). The soft state's observable weight
disappears. Conversely, a positive gap at each \(n\) can collapse with
nonvanishing weight. These examples identify the estimates a continuum
argument must actually carry.

For Newton the analogous surviving object is a calibrated record and its
information about the comparison. No spectral measure is supplied by
classical mechanics for this purpose. The useful transfer is the proof
discipline: compatible observable limits, nonvanishing response, and a
bound uniform over admitted refinements. Any stronger equivalence
requires a constructed mathematical map.

\hypertarget{the-mechanism-worth-developing-across-the-two-main-goals}{%
\subsection{6. The mechanism worth developing across the two main
goals}\label{the-mechanism-worth-developing-across-the-two-main-goals}}

The \href{action-floor-yang-mills-gap.md}{action-floor note G08}
supplies one concrete bridge. A transverse oscillator of frequency
\(\Omega(x)\) obeys

\[\frac{p_y^2}{2M}+\frac12M\Omega(x)^2y^2
\ \ge\ \frac12\hbar\Omega(x).\]

When \(\Omega(x)\) grows along a classically escaping valley, this
all-state operator bound can confine that valley. Its force comes from
the action scale together with the non-abelian commutator potential. The
solved matrix models establish the mechanism in a specified Hamiltonian;
extension to four-dimensional gauge theory requires uniform control of
the additional modes and renormalization.

The pion comparison asks which flat directions an interaction is allowed
to lift. A broken exact global symmetry protects Goldstone directions.
Pure Yang--Mills has no corresponding axial flavour symmetry. The
absence of that protection makes lifting possible; the positive bound
still has to be proved. Finite quark mass then supplies a controlled way
to lift the protected direction, which is why both pion cases are useful
as orientation.

The candidate unifying question is therefore: \textbf{which structures
enforce an irreducible cost, which protect a soft direction, and which
survive the removal of the regulator?} Quantum kinematics supplies an
action unit to both QCD theories. Their different spectra depend on
dynamics and symmetry. A successful Newton necessity argument must
explain the action structure at the preceding level.

\hypertarget{work-toward-the-final-paper}{%
\subsection{7. Work toward the final
paper}\label{work-toward-the-final-paper}}

The intended paper can be called \emph{What survives refinement: action,
symmetry and the continuum}. Its architecture is: Galileo's comparison
and the three regulators; the pion benchmark with and without explicit
breaking; physical spectral limits and reconstruction; the established
Newton record bounds; the valley-lifting bridge; and the two remaining
proof obligations. Historical claims stay tied to the held source
companions and the sibling \texttt{newtonlean} work. The present Planck
paper becomes a component of this synthesis. This working architecture
does not designate either open goal as solved.

The research order is deliberately bounded.

\begin{enumerate}
\def\labelenumi{\arabic{enumi}.}
\tightlist
\item
  \textbf{Next constructive step: the \(SU(3)\) small-volume bridge.}
  Develop the H3 estimate for the zero-mode Hamiltonian plus the even
  torus valley potential in the
  \href{weak-coupling-feshbach-reduction.md}{Feshbach note}. Seek an
  explicit lower bound \[E_{\rm small}(L)\ge
  (d_3-Cg(L)^{2/3})g(L)^{2/3}\frac{\hbar c}{L},\qquad d_3>0,\] with an
  explicit admissible coupling interval. State the operator, gauge
  sector and cutoff dependence. A confining potential or a positive
  ground energy alone does not establish the excitation gap; the first
  excited energy and the ground energy need comparison. Stop at one
  proved estimate, or the precise uncontrolled term in that estimate.
  Fixed-cutoff H3 would be a result, with its boundary retained; it
  would not close H1/H2 or T4 of the full mass-gap map.
\item
  \textbf{Newton necessity: construct the missing physical premise.} Use
  one explicit model of records and their composition, with calibrated
  lengths, impulses and time. Seek a refinement-independent positive
  cost from dynamical or optical premises that have their own
  justification. The deliverable is one implication with its full
  admissible class, or an explicit identification of the additional
  premise needed. A result starting from the canonical commutator or M3
  belongs to the existing conditional branch. For the optical route,
  observable relative phase and a common action unit across probes must
  be established separately. Attainment of the existing disturbance
  bounds remains useful, but secondary to this task.
\item
  \textbf{Use the pion as a bounded test of a proposed mechanism.} Carry
  a non-singlet Ward identity and its nonzero residue through the chosen
  regulator, or state these as explicit inputs. Recover the zero-mass
  case and its small-mass lifting. This is one comparison section; a
  full fermionic-QCD construction is outside the active queue.
\item
  \textbf{Return the successful bridge to the field theory.} Determine
  whether the small-volume estimate survives eliminating nonzero modes
  with constants uniform in \(a\), and then whether it supplies the
  order-one-coupling blocking estimate missing from H1. Preserve the
  chosen vacuum, physical scale and observable normalization.
  Reconstruction and nontriviality stay separate obligations; their
  order may interact with the estimates rather than follow a single
  irreversible chain.
\item
  \textbf{Integrate proved results into the final synthesis.} Replace
  the corresponding conditional statements only when their hypotheses
  are discharged. Complete the historical reading obligations in the
  Planck paper before submission. Keep the pion's role explanatory
  unless it actually changes one of the two proofs.
\end{enumerate}

\hypertarget{consequence-for-state}{%
\subsection{8. Consequence for STATE}\label{consequence-for-state}}

The destination becomes a joint paper on survival under refinement, with
the Newton action necessity and pure \(SU(3)\) continuum mass gap as the
two main goals. Pion physics, at zero and nonzero quark mass, is the
symmetry and limit-order benchmark requested by the user. The next
bounded constructive task is the H3 small-volume bridge; the Newton
necessity task retains the independent-premise requirement. The three
review batches of the existing Planck paper are complete; attainment and
its submission obligations remain open. This note is the live plan, and
adds no claim of a four-dimensional construction or an independent
derivation of \(h>0\).
\end{document}
